%%
\documentclass[manuscript,screen,review]{acmart}

%% ===================== 中文支持 =====================
%% scheme=plain: 只接管中文字体与断行,不改写 acmart 的章节名。
%% fontset=windows: 本机宋体/黑体; Overleaf/Linux 可改为 fontset=fandol。
\usepackage[UTF8,scheme=plain,fontset=windows]{ctex}

\usepackage{algorithm}
\usepackage{algorithmicx}
\usepackage{algpseudocode}
\usepackage{pifont}
\usepackage{makecell}
\providecommand{\checkmark}{\ding{51}}
%% B13 删去 \cmark/\xmark/\omark:定义后全稿 0 处引用(表格未采用勾叉体例)。原值见 git。
\usepackage{threeparttable}
\usepackage{tikz}
\usetikzlibrary{positioning,arrows.meta,calc,shapes,shadows}
\usepackage{xcolor}
\usepackage{amsmath}
\usepackage{bm}
\DeclareMathOperator*{\argmin}{arg\,min}
\DeclareMathOperator*{\argmax}{arg\,max}
\usepackage{multirow}
\usepackage{enumitem}
\usepackage{booktabs}
\usepackage{graphicx}
\setlength{\headheight}{20.48303pt}

%% scheme=plain 不改浮动体名称,与 paper01 中文稿对齐。
\renewcommand{\figurename}{图}
\renewcommand{\tablename}{表}
\floatname{algorithm}{算法}
\renewcommand{\proofname}{证明}

%% 机制缩写:主语是闭包,不是「修复框架」
\providecommand{\STRC}{\textsc{STRC}}
%% B13 删去 \TBD:全稿 0 处引用,待填项一律用 <待填:...> 明文标记,便于自检抓取。

%% =====================================================================
%% 主结果数字的单一来源。
%%
%% 全部由 STRC/tools/paper_numbers.py 从 STRC/experiments/expanded/*.csv 打印,
%% 数据由 `py -m tools.expand_batch` 生成(默认 5 算例 x 10 种子)。
%% 正文一律用宏,不写字面量;改数只改这里,并同步跑一遍上述两条命令核对。
%%
%% !! 2026-08-16 由 3x5=15 对扩到 5x10=50 对。旧稿所有 /15 结果一律作废。
%%
%% !! 2026-08-16(二次):E2b 由 24/50 变为 50/50。这不是换了算例,而是修掉了
%% !! 一处实现与定义的错配。释放集是一个**预约**集合,而 replay_reroute 原先
%% !! 按**工序**判定是否重路(两段中任一段进闭包就整道重规划)。于是当只有满载
%% !! 段进闭包时,同工序的空载段也被改写,而它的 task 并不在释放集里,遂被记为
%% !! 外侧漂移。26 个失败格上的 84 条漂移预约无一例外都是这种同工序空载段,
%% !! 其中 71 条(85%)在 t_now 之前就已执行完毕——改写它同时违反假设 A2。
%% !! 改为按段判定后漂移归零。据此 prop:outside 的前提 (iii) 才真正被满足,
%% !! 此前的 24/50 测的是实现缺陷,不是边集的物理覆盖度。
%% =====================================================================
\newcommand{\NInst}{5}
\newcommand{\NSeeds}{10}
\newcommand{\NPairs}{50}
%% E1/E3:两项都是满分,写成 \NPairs/\NPairs 以免手抄漏改
\newcommand{\PassMiss}{\NPairs/\NPairs}      % T_impact=0 且种子非空且原解不可行
\newcommand{\FeasRone}{0/\NPairs}            % R1 恢复可行的格数
\newcommand{\FeasRtwo}{\NPairs/\NPairs}      % R2 恢复可行的格数
\newcommand{\PassClosed}{\NPairs/\NPairs}    % E2a 结构闭合(泄漏为 0)
\newcommand{\PassInvar}{\NPairs/\NPairs}     % E2b 外侧字段逐条不变
%% E1 主批次 Cl/|R| 的逐算例极值,出处 tab:e1 表体(0.59/0.57/0.59/0.57/0.56)。
%% 此前该区间在正文三处各写一个值(0.56--0.59 / 0.54--0.57 / 0.54--0.58),即 CONFLICT-3;
%% 现统一由本宏出数。**不要**与下列同值但不同源的量混用:
%%   \ScaleCloHi(0.56) 来自规模扫描 tab:scale 的 k=4 格,是另一批次;
%%   \PubEoneFracLo/Hi 来自外部布局批次,账本独立。
\newcommand{\EoneCloLo}{0.56}
\newcommand{\EoneCloHi}{0.59}
%% 耗时(E3 批次,第 1 级改路、无扩域)
\newcommand{\MsRepairMed}{3.2}
\newcommand{\MsRepairMax}{12.8}
%% 稳定性:E5 上受 A2 约束后两臂改动比例已同区间;主对照改为 E7 的 R2 vs RA。
\newcommand{\ResFracRtwo}{50\%--66\%}
%% E5 批次:R0+ 已走固定前缀解码。降级后道改走同车续路后,18 个预算点越界均为 0。
%% 出处 expanded/e5_cross_curve.csv,`py -m tools.expand_batch --only-e5`。
\newcommand{\RzeroPastLo}{0}
%% ---------------------------------------------------------------------
%% E5 两批合读。pub_batch 直接 import expand_batch 的 _run_e5,
%% t_now_frac=0.35、预算 {0.2,1,2}s、pop=40、hot=True 逐项同一,故可合池;
%% 两批只差布局出处(自建 / 外部转录),这正是合池的意义。
%% 出处 expanded/e5_cross_curve.csv + pub_layouts/e5_cross_curve.csv,
%% 复算脚本 tools/_p17_e5_stats.py。
\newcommand{\EfiveInst}{4}          % congested_8x4x4, example_3x3x2, LyuL2_4m, LyuL6_8m
\newcommand{\EfiveCells}{12}        % 独立算例--随机种子格
\newcommand{\EfivePts}{36}          % 预算点总数 = 12 格 x 3 档
\newcommand{\EfiveWLT}{35/1/0}      % 逐预算点 R0+ 胜/持平/负
\newcommand{\EfiveWinMin}{11}       % 最小预算档上 R0+ 严格更优的格数
\newcommand{\EfiveTieMin}{1}        % 同,持平的格数(LyuL6_8m, 种子 2024)
%% ---------------------------------------------------------------------
%% 离散度口径(P17)。均值表/图一律配一个四分位距,免得只读到均值。
%% 出处同上脚本;E1 为 expanded/e1_miss.csv 逐算例 10 种子。
\newcommand{\EoneFracIqrLo}{0.024}  % 逐算例 Cl/|R| 的 IQR 最小值
\newcommand{\EoneFracIqrHi}{0.049}  % 同,最大值
\newcommand{\EoneFracMed}{0.571}    % 50 对合池中位
%% 合池 IQR 实测 0.0372,与 §5 的 Hodges--Lehmann 估计 0.037 同值\emph{不同源};
%% 为免两个不相干的量在正文里撞成一个结果,此处不出宏、图注也不报该值,
%% 只报逐算例 IQR 区间与合池极差,已足够说明离散度。
\newcommand{\EoneFracMin}{0.500}    % 50 对合池极差下端
\newcommand{\EoneFracMax}{0.655}    % 同,上端
%% phi 扫描(tab:scale / fig:scale):耗时极差与 Cmax 逐档 IQR
\newcommand{\ScaleMsMin}{2.2}       % 全扫描 strc_wall_ms 极差下端
\newcommand{\ScaleMsMax}{23.8}      % 同,上端
\newcommand{\ScaleCmaxIqrLo}{8}     % 逐档 R0+ Cmax 的 IQR 最小值
\newcommand{\ScaleCmaxIqrHi}{76}    % 同,最大值
%% 规模扫描的加速比区间(跨两算例的逐格极值)
\newcommand{\SpeedLo}{87}
\newcommand{\SpeedHi}{3051}
%% E6:A 类上闭包严格大于工序依赖图释放集的格数(其余格两者相等,不是包含失败)
\newcommand{\StrictlyLarger}{87/100}
%% E6:闭包占 t_now 之后活预约的比例(走廊阻断)。这个数偏大,
%% 是「不作最小性主张」那条注记的实测依据,不要美化。
\newcommand{\ClosureAliveFrac}{92\%}
%% =====================================================================
%% 外部来源布局批次(STRC/experiments/pub_layouts/)。**独立账本**,不并入上面的
%% \NPairs 对——把它混进主批次会让全部 /\NPairs 结果改口径。两批算例逐参数同口径,
%% 只差布局来源。数出处:py -m tools.pub_batch 后 py -m tools.paper_numbers。
%% 边权是本文补齐的等权值,故这批数**不可**与 Lyu 或 van Os 的参照值比较。
%% =====================================================================
\newcommand{\NPubInst}{5}
\newcommand{\NPubPairs}{50}
\newcommand{\PubPassMiss}{\NPubPairs/\NPubPairs}
\newcommand{\PubFeasRone}{0/\NPubPairs}
\newcommand{\PubFeasRtwo}{\NPubPairs/\NPubPairs}
\newcommand{\PubPassClosed}{\NPubPairs/\NPubPairs}
\newcommand{\PubPassInvar}{\NPubPairs/\NPubPairs}
\newcommand{\PubEoneFracLo}{0.438}
\newcommand{\PubEoneFracHi}{0.547}
%% E4 同参数对照。**这里的口径要格外小心,曾写错过一次。**
%% \NPubInst 个外部算例只落在 \NPubNets 张不同路网上(4x4 去一条边:LyuL2/L3;
%% 5x5 去两条边:LyuL4/L5/L6),差别在机器台数与落位。于是:
%%   \PubCloSpread 这个极差的两个端点(\PubCloHi=LyuL2 4 机、\PubCloLo=LyuL3 5 机)
%%   **共享同一张路网**,所以它测的是机器落位,**不是**路网几何。
%%   不要再写成"闭包占比随布局显著变化"——数据不支持这句。
%% 正确的用法是:同一路网上换落位即得 \PubCloSpread,而 LU 最小割与漏斗占比全程恒定
%% (两张路网都把装卸站放在网格角点,四邻接下角点恰两条出边),故静态割被干净否证。
%% 注意别把这句写成"公开布局一律放在角点":工具集另附的 Liu 那族放在网格边中点(出度 3),
%% 只是因允许对角移动而未收入。准确的说法是该指标只取到极少数由摆位决定的离散值。
\newcommand{\NPubNets}{2}
\newcommand{\PubCloLo}{0.305}
\newcommand{\PubCloHi}{0.512}
\newcommand{\PubCloSpread}{0.207}
\newcommand{\PubCloSpreadBig}{0.075}
\newcommand{\SelfCloLo}{0.446}
\newcommand{\SelfCloHi}{0.497}
\newcommand{\SelfCloSpread}{0.051}
\newcommand{\PubEfiveNoCross}{17/18}
%% =====================================================================
%% 便宜对照臂(E7)与算例规模扫描(E8)。两批共用同一组臂定义:
%%   RS  全局右移。不改路径/指派/序,把未完成的一切统一后推到阻断窗之后。
%%       冻结判据与 R2 相同(t_end <= t_now 不可动),故按 A2 可采纳。
%%   RD  原染色体重解码。保持机器指派与加工顺序,走与 R0+ 同一套固定前缀解码。
%%       降级后道改走同车续路后,A2 越界为 \RedecodeBad。
%%   RA  不划边界:释放 t_now 之后的全部预约,再走与 R2 完全相同的改路重放。
%%       它与 R2 共用引擎、共用冻结判据,唯一差别是释放集取了平凡上界,
%%       故 R2--RA 之差可全部归因于闭包这个对象本身。
%% 出处:STRC/experiments/cheap_baselines.csv (py -m tools.cheap_baselines)
%%      STRC/experiments/scale_curve.csv     (py -m tools.scale_curve)
%% 协议与 E1--E3 一致(繁忙走廊单窗阻断、t_now=0.35 Cmax、R2 不开扩域),
%% 故 tab:cheap 可与 tab:e3 并读;但**不可**与 tab:e5 并读——那批只有两个算例。
%% =====================================================================
\newcommand{\MsShift}{1.42}          % RS 耗时中位/ms(\NPairs 对)
\newcommand{\MsRedecode}{4.16}       % RD 耗时中位/ms(同上)
\newcommand{\MsAllFuture}{2.76}      % RA 耗时中位/ms(同上)
\newcommand{\MsCheapRtwo}{2.74}      % 同批 R2 的耗时中位,供四臂同口径比较
\newcommand{\RedecodeRatio}{1.5}     % MsRedecode / MsCheapRtwo
\newcommand{\ShiftBad}{0/\NPairs}    % RS 改写已完成预约的格数
\newcommand{\RedecodeBad}{0/\NPairs}
\newcommand{\WinVsShift}{27/16/7}    % R2 对 RS 的 Cmax 赢/平/输
\newcommand{\WinVsRedecode}{14/14/22}
\newcommand{\WinVsAllFuture}{4/46/0}
\newcommand{\ResFracCheapRtwo}{0.540}
\newcommand{\ResFracShift}{0.627}
\newcommand{\ResFracRedecode}{0.618}
\newcommand{\ResFracAllFuture}{0.604}
%% R2 对 RA 的逐格稳定性对比(改动比例;两臂都可行的 \NPairs 格全部纳入)。
\newcommand{\RAStabWTL}{31/18/1}     % R2 更稳/持平/更差
\newcommand{\RAStabP}{8.8\times10^{-7}}  % Wilcoxon 符号秩,双侧
\newcommand{\RAStabMean}{0.064}      % 改动比例平均降幅
%% 下面两个是剔除 6 个「两法恒等」格(R2 释放集 = 全部未结束占用)后的同一结果,
%% 由 STRC/tools/audit_checks.py 自 cheap_baselines.csv 直算;剔除使降幅变大。
\newcommand{\RAStabWTLNet}{31/12/1} % 剔除恒等格后的 R2 更稳/持平/更差(n=44)
\newcommand{\RAStabMeanNet}{0.073}  % 同,平均降幅
\newcommand{\RAStabMed}{0.029}       % 同,中位降幅
\newcommand{\RACmaxRatio}{0.999}     % Cmax 比 R2/RA 均值
\newcommand{\RACmaxLo}{0.970}        % 同,最小值(R2 最占优的那格)
\newcommand{\ClosureShareLo}{0.80}   % 闭包/全部未来预约:最小值
\newcommand{\ClosureShareMean}{0.92} % 同,均值(与 \ClosureAliveFrac 同源)
%% 算例规模梯子:工件 2k、机器 k、车 k,k=4..10;拥堵档、H、F、Tt/Tp 与两处割集同口径,
%% 只有规模变(生成器 clbs/tools/gen_instances.py 的标定见其文档字符串)。
\newcommand{\ScaleKs}{7}
\newcommand{\ScalePairs}{70}
\newcommand{\ScaleJobsHi}{20}
\newcommand{\ScaleMachHi}{10}
\newcommand{\ScaleCloLo}{0.48}       % 闭包占全表比例:k=10
\newcommand{\ScaleCloHi}{0.56}       % 同,k=4
%% 下面三个与上面两个是同一批数据的两种分母,由 tools/audit_checks.py 自 scale_curve.csv
%% 的 n_res / n_live / closure 三列直算(不是反推)。Cl/|R| 的可达上界是 |R^o|/|R|,
%% 故回答「占比会不会涨到接近 1」必须用 Cl/|R^o| 这一口径。
\newcommand{\ScaleAliveShare}{0.61}  % 均 |R^o|/|R|,七档均值 0.611
\newcommand{\ScaleCloAliveLo}{0.78}  % 闭包占尚未结束占用之比,档均值最低:k=10 的 0.782
\newcommand{\ScaleCloAliveHi}{0.91}  % 同,最高:k=4 的 0.905
\newcommand{\ScaleMsHi}{10.4}        % R2 耗时中位:k=10
\newcommand{\ScaleRatioLo}{1.5}      % RD/R2 耗时比:k=4
\newcommand{\ScaleRatioHi}{3.3}      % 同,k=10
\newcommand{\ScaleFeasFirst}{69/\ScalePairs}   % 第 1 级改路即可行的格数
\newcommand{\ScaleFeasAll}{\ScalePairs/\ScalePairs}
%% E7 事后规则:释放占比 vs 稳定性降幅。阈值在 cheap_baselines.csv 本样本上选定,
%% 不作外推定理。出处:py -m tools.paper_numbers(worth_section)/fig_strc_worth_CN.py
\newcommand{\WorthThresh}{0.93}
\newcommand{\WorthLowN}{27}
\newcommand{\WorthLowMean}{0.094}
\newcommand{\WorthHighN}{23}
\newcommand{\WorthHighMean}{0.029}
\newcommand{\WorthVanish}{0.97}
%% 边集差分试探。出处:STRC/experiments/edge_probe.csv (py -m tools.edge_probe --full)
\newcommand{\ProbeLeakPairs}{0/50}
\newcommand{\ProbeLeaks}{0}
%% E6 修复维(第 1 级、不开扩域)。出处:expanded/e6_types.csv 的 R2_feasible 列
\newcommand{\FeasESixBlock}{\NPairs/\NPairs}
\newcommand{\FeasSlow}{\NPairs/\NPairs}
\newcommand{\FeasESixAgv}{\NPairs/\NPairs}
\newcommand{\FeasESixRa}{\NPairs/\NPairs}
%% 非均匀边权 E1/E3。出处:STRC/experiments/weight_sweep.csv (py -m tools.weight_sweep)
\newcommand{\WeightLogRtwo}{\NPairs/\NPairs}
\newcommand{\WeightLuRtwo}{\NPairs/\NPairs}
\newcommand{\WeightSplitRtwo}{49/\NPairs}

\usepackage{amsthm}
\newtheorem{assumption}{假设}
\newtheorem{definition}{定义}
\newtheorem{problem}{问题}
\newtheorem{remark}{注}
\newtheorem{prop}{命题}
\newtheorem{lem}{引理}

\AtBeginDocument{%
  \providecommand\BibTeX{{%
    Bib\TeX}}}

\setcopyright{acmlicensed}
\copyrightyear{2026}
\acmYear{2026}
\acmDOI{XXXXXXX.XXXXXXX}

%%\acmSubmissionID{123-A56-BU3}

\begin{document}

%% =====================================================================
%% 标题:独立小论文用「工序依赖图无法刻画」钩子;大论文章题另用「闭包上的有界修复」
%% B11 定名:题名与正文统一用「工序依赖图」(CONFLICT-4);删「级联」二字——该词全稿
%% 仅题目一处,正文与贡献、结论一律作「失效」,删后逐字一致且全标题回到 29 字(≤30)
%% =====================================================================
\title[工序依赖图无法刻画的走廊占用失效]{时空预约闭包：柔性作业车间中工序依赖图无法刻画的走廊占用失效}

\author{Qi Yu}
\email{230240012@seu.edu.cn}
\orcid{0009-0008-7843-6345}
\affiliation{%
	\department{School of Cyber Science and Engineering}
	\department{Key Laboratory of Computer Network and Information Integration}
	\institution{Southeast University}
	\city{Nanjing}
	\state{Jiangsu}
	\country{PR China}
}

\author{Wei Li}
\orcid{0009-0001-0949-329X}
\email{xchlw@seu.edu.cn}
\authornote{corresponding author}
\affiliation{%
	\department{School of Computer Science and Engineering}
	\department{Key Laboratory of Computer Network and Information Integration}
	\institution{Southeast University}
	\city{Nanjing}
	\state{Jiangsu}
	\country{PR China}
}

\renewcommand{\shortauthors}{Yu and Li}

%% =====================================================================
%% 摘要
%% 骨架:背景缺口 → 主张(边界在预约图) → 对象 STRC → 用法(有界修复) →
%%       与全局重解分工(响应时间一侧) → 实验钩子(不把创新挂在「又一个修复算法」上)
%% =====================================================================
\begin{abstract}
在同时用工序依赖刻画加工先后、又用时间窗预约记录走廊占用的柔性作业车间调度中，
扰动后划定重调度范围的现成做法仍是先找出受影响工序、再连同其后续一并重排。
但走廊在某段时窗内不可通行或通行变慢时，没有任何工序因此无法执行，
该做法给出空范围、判为无需重调度；而落在受扰时段内的走廊占用已经失效，原排程不可行，
等待还会沿让行关系传给后续车辆。缺口不在检测，而在工序依赖图
不含走廊占用之间的让行关系。

本文把重调度范围定义在走廊占用及其依赖关系上，提出时空预约闭包(\STRC)，
其结果集合称为预约影响集：以占用时段与扰动时窗重叠的走廊预约为起点，沿同走廊让行、
同一车辆的后续占用、同一工件与同一机器下一道工序的占用四类边取传递闭包；
集内占用允许改写，集外沿用原计划。依赖边一经给定，一条占用是否属于该集即唯一确定，
不随改路算法的实现而变。恢复分两级：先在集内为原车辆改路，仍不可行时再按既定顺序
扩大允许改写的范围。

实验在 \NInst\ 个算例 \(\times\) \NSeeds\ 个随机种子共 \NPairs\ 对上进行。
走廊阻断时，按工序依赖图划出的范围在 \PassMiss\ 对上均为空，而原排程均已不可行。
改路程序不变、只替换划界对象时，恢复可行率由 \FeasRone\ 变为 \FeasRtwo。
预约影响集对外封闭（集内无指向集外的依赖边）与集外占用的走廊、时段和车辆逐条不变
分别在 \PassClosed\ 与 \PassInvar\ 对上成立。若把允许改写的集合换成
「决策时刻之后尚未结束的全部占用」、其余不变，耗时与完工时间几乎不变，
改动比例则由 \ResFracCheapRtwo\ 升至 \ResFracAllFuture（逐格 \RAStabWTL,
Wilcoxon 符号秩双侧 \(p=\RAStabP\)）。与不改写已结束占用的热启动全局重调度相比，
有界修复耗时低两到三个数量级，但完工时间在全部预算档上均更差，且放宽预算后不出现反超。

主效应来自划界对象的替换：由工序依赖图换为预约影响集后可行性得以恢复，改路程序未变。
次效应来自预约影响集相对「尚未结束的全部占用」的缩减，即改动更少；但前者已覆盖后者的约
\ClosureAliveFrac，两者接近时该收益趋于消失。
集内改路、集外保持与失败后扩大范围沿用既有做法，本文的贡献在于划界对象。
\end{abstract}

\keywords{柔性作业车间调度, AGV 无冲突路由, 时空预约表, 时空预约闭包, 走廊阻断}

%% CCS。投 ACM 时须用官方工具重生成 CCSXML 块;此处先给 \ccsdesc 以消除
%% acmart 对超过两页稿件的 "no CCS concepts" 警告。非 ACM 场合整段可删。
\ccsdesc[500]{Theory of computation~Scheduling algorithms}
\ccsdesc[300]{Applied computing~Industry and manufacturing}
\ccsdesc[300]{Computing methodologies~Planning and scheduling}

\maketitle

%% =====================================================================
\section{引言}
\label{sec:introduction}
%% =====================================================================

\subsection{背景与动机}

考虑运输的柔性作业车间调度长期建立在一个简化假设之上：工件在两台机器之间的移动时间是一个常数，可从按位置对索引的矩阵中读出。
该假设默认车辆互不争用同一路段，行驶时间因此只由两点之间的距离决定，而不是路由决策的结果。
一旦有限车队共享稀疏走廊，并以时间窗预约被争用的路段，行驶时间就不再是固定数据，而成为两个决策的结果：车辆的行驶路径，以及各车辆占用该路段的先后次序。
一辆车的让行会推迟另一辆车的到达，进而推迟后续工序开工及其运输占用。
此时一份可行排程须同时确定工序指派、加工顺序，以及一张记录走廊占用的预约表；扰动既可能使工序无法执行，也可能只使走廊占用失效而不改变任何工序指派。

在开工前一次性完成排程、并把机器指派与无冲突路径一同决定的这类研究已经表明：
当工序可选多台机器时，车辆在走廊上的争用会改变哪台机器更优——通往加工更快的机器的路段若正被占用，加工较慢但更易到达的机器反而更好。
因此，评价候选排程须用时间窗路由计入等待，搜索时还须查询预约表上的冲突，
而不能继续使用常数行驶时间~\cite{vanos2026exact,li2026framework,sun2023integrated}。
上述文献回答的是：在信息完备、尚未开工时，如何排出一份高质量的无冲突排程。同一车间在执行到时刻 \(t_{\mathrm{now}}\) 并发生扰动之后，
问题随之变为：失效落在哪些走廊占用上，改写哪些占用才足以恢复可行，以及能否在不重新启动种群搜索的前提下于毫秒量级内完成恢复。

\subsection{工序依赖图无法刻画的缺口}

%% 措辞收窄(2026-08-16):原句「局部重调度文献常用任务依赖图界定影响域」是一句
%% 针对整支文献的全称否定,一个反例即可推翻——铁路 alternative graph 与 MAPF 都不
%% 在工序依赖图上划界。收窄为「考虑运输的柔性作业车间动态调度这一支」,见 2.5。
在考虑运输的柔性作业车间动态调度研究中，重调度范围通常沿工序依赖图划定。
具体做法是：以因设备故障而无法按原计划执行的工序为起点，把同一工件的后续工序与同一机器上排在其后的工序纳入范围，
或在依赖图上从起点扩展至多 \(\theta\) 步后停止；随后只对范围内的工序重调度。
两篇带 AGV 运输、并以机器故障或插单为扰动的动态调度研究即按此划界——恢复范围分别取「故障后尚未加工的工序」\cite{zhang2023energy}，
与「按扰动影响分为保留、继续、重构三类后的工序与配送任务」\cite{tang2026dynamic}（后者为批量流混合流水车间，扰动含车辆故障，但模型中没有预约表）。
这一做法可上溯至受影响工序重调度（AOR）\cite{zhang2023energy,tang2026dynamic}。
后文以按工序依赖图划界的对照代表这一做法（第~\ref{sec:rw_aor}~节）。

这一划界在扰动本身就是工序事件时往往合理——例如机器故障使若干工序无法按原指派执行。
然而，当排程用时间窗记录走廊占用时，还存在一类扰动：它使某条走廊在某段时间内不可通行，而不改任何工序指派。典型例子是走廊阻断：
走廊在 \([t_a,t_b)\) 时段禁止通行。阻断不修改机器指派或加工顺序，工序依赖图上因此没有任何被直接波及的工序，按该方法划出的重调度范围为空，并据此判为无需重调度。
于是工序依赖图无法刻画这些必须恢复的走廊占用：原计划中与阻断时段重叠的占用已与禁行冲突，不再可行；一辆车因此等待或改道，又会沿让行关系阻塞后续车辆，失效由此在预约表上传播。

这里的「无法刻画」指工序依赖图按定义不含走廊占用之间的让行关系，因而无法表达占用失效在车辆之间的传播。
若运输时间仍取为按位置对索引的常数矩阵，走廊不是可争用的资源，预约表也不存在，这类失效在模型中无从表达；
因此该缺口只存在于用时间窗记录走廊占用的设定中。

据此，本文区分两类扰动。
\begin{itemize}[leftmargin=1.5em,itemsep=0.25em]
\item \textbf{A 类}（作用于工序依赖图）：机械臂故障、加工时间显著偏差、插单，以及\emph{车辆故障}。
这类扰动使部分工序无法按原计划执行，划定重调度范围时的起点工序集非空；
按工序依赖图划界与按预约影响集划界都能给出非空范围，因而不是本文用来检验划界错位的主要对象。
车辆故障归入 A 类的理由见第~\ref{sec:disturbance_classes}~节：它虽不改工序指派，
却使该车承运的工序不可执行，经「车辆到所承运工序」的对应关系，工序依赖图上仍有非空的起点工序集。
\item \textbf{B 类}（只作用于预约表）：走廊阻断、走廊通行变慢，以及局部通行权临时取消等。
这类扰动不使任何工序无法执行，工序依赖图上的起点工序集为空；
与受扰时段重叠的走廊占用非空，以其为起点纳入重调度范围所得的预约影响集也非空。
\end{itemize}
B 类扰动构成对「按工序依赖图划界」这一做法的问题层面反例：模型同时用工序依赖图刻画加工先后、
用时间窗记录走廊占用时，若重调度范围仍只划在工序依赖图上，便无法表达走廊阻断这类扰动。
上引两篇带 AGV 的动态调度研究不考虑车辆之间的走廊争用（第~\ref{sec:rw_position}~节），
工序依赖图是其唯一可用的划界对象；这一错位只在两套对象同时出现时发生（第~\ref{sec:rw_claim}~节）。
图~\ref{fig:motivate} 对照这一缺口：左侧工序依赖图上没有作为起点的工序，因而判为无需重调度；
右侧预约表上，与阻断时段重叠的占用成为起点，并沿让行与接续关系形成预约影响集。

\begin{figure}[t]
\centering
\includegraphics[width=\linewidth]{figures/fig_strc_motivate_CN}
\caption{动机示意。左：走廊阻断不使任何工序或机器无法按原计划执行，工序依赖图上没有作为起点的工序。
右：同一扰动下，与阻断时段重叠的走廊占用成为起点，并沿让行与接续关系形成预约影响集。}
\Description{并排两幅图。左幅是工序依赖图，走廊阻断没有使任何工序节点无法执行，
作为起点的工序集为空。右幅是同一时刻的走廊预约表，与阻断时段重叠的占用成为起点，
这些占用沿让行边与接续边向外连成一片传播区域。}
\label{fig:motivate}
\end{figure}

\subsection{本文主张}

本文有三点主张。

第一，当柔性作业车间用时间窗预约保证车辆无冲突通行时，扰动后的重调度范围应定义在由走廊占用及其让行关系构成的预约依赖图上，而不是只定义在工序依赖图上。

第二，划界方法是时空预约闭包（\STRC, Spatiotemporal Reservation Closure），其结果集合称为预约影响集：
以与扰动时窗重叠、因而失效的走廊占用为起点，沿让行、同一车辆的后续占用、同一工件下一道工序的占用、同一机器上下一道工序的占用，
把直接或间接必须作废、重写的走廊占用全部纳入重调度范围。
这些依赖关系一经给定，影响集以外的占用不再是集内占用沿上述关系的下一步，
因而「改写哪些占用」对应一个可以检验的集合，而不是人为设定的搜索半径。

第三，得到影响集之后，在集内为原车辆改路、集外保持原占用；若仍得不到可行排程，
再按同一工件、同一车辆、直至全部未完成占用的顺序，扩大允许改写的范围。
A 类故障上另允许换车与改派，仍不改加工顺序。
这些步骤沿用既有做法，用于落实影响集并恢复可行排程。

上述三点共用一个立场：\textbf{预约影响集是一次测量，不是一个决策。}
本文不用占用间的让行关系引导搜索以更换指派或路径，因为那属于改进排程决策：
消除一处走廊等待后，另一处约束往往成为新的瓶颈，完工时间上的收益因此容易被抵消。
本文用同一组让行关系划定哪些占用已经失效：依赖关系给定后，
一条占用是否属于影响集，取决于它能否从起点沿这些关系到达，结果唯一确定，
不随之后哪一条约束成为新的瓶颈而改变。

纳入哪些依赖边是建模选择；边集选定后，影响集对外封闭（集内占用没有指向集外的依赖边）
即为定义的推论（命题~\ref{prop:struct}）。
影响集不一定是恢复可行所需的最小占用集合，因为最小集合依赖于允许的修复动作，
而影响集的定义不绑定于任何一种修复方法（注~\ref{rem:minimality}）。

相对以原排程为起点、在限定时间内重新搜索的全局重调度，预约影响集回答的是另一类问题：
在必须很快给出可执行方案时，能否恢复可行，并把改写限制在影响集以内。
在命题~\ref{prop:outside} 的条件下，修复后影响集以外每条走廊占用的走廊、时段与车辆均与原计划一致。
相对不改写决策时刻已经结束的占用的全局重调度，有界修复得到的完工时间 \(C_{\max}\) 更差，
两者的改动比例处于同一量级；有界修复的长处在于算法耗时短得多（毫秒量级对秒量级）。
算法耗时与改动比例须分开解释，第~\ref{sec:cheap}~节对二者分别检验：
耗时短来自对影响集内各段运输各做一次路径规划、不重新启动种群搜索，而不来自影响集本身；
影响集是否减少了改动，无法从上述全局重调度对照中得出（两者比例已经接近），
须另设对照：改路程序相同，唯一差别是允许改写的集合取影响集还是尚未结束的全部占用。
在所测的 \(\EfiveCells\) 个算例--随机种子格上，全局重调度在各搜索时限下得到的 \(C_{\max}\) 均更小。
按工序依赖图划界时，走廊阻断因找不到受影响工序而不会触发局部重调度，原排程保持不可行；
若改为触发全局重调度，则可以恢复可行，但算法耗时为秒级。
因此，本文关注的是须立即恢复可行时走廊占用上失效范围的识别及相应的局部恢复。

\subsection{贡献}
\label{sec:contributions}

本文的贡献如下，后文分别记为 C1、C2、C3。

\begin{enumerate}[label=\textbf{C\arabic*.},leftmargin=2.2em,itemsep=0.4em]
\item \textbf{指出工序依赖图无法刻画的一类失效。}
当排程同时用工序依赖刻画加工先后、又用时间窗记录走廊占用时，
现成的重调度范围仍按受影响工序划定。
走廊阻断、走廊通行变慢等扰动不改任何工序指派，
按该方法划出的范围为空，并据此判为无需重调度；
但与受扰时段重叠的走廊占用已经失效，原排程不可行，
等待还会沿让行关系传给后续车辆。
其根源在于工序依赖图不含占用之间的让行关系。

\item \textbf{把重调度范围定义为预约影响集。}
以与扰动时窗重叠的走廊占用为起点，
沿让行、同一车辆的后续占用、同一工件与同一机器上下一道工序的占用，
把直接或间接必须作废、重写的占用纳入重调度范围；这一划界方法称为时空预约闭包（\STRC）。
依赖关系一经给定，一条占用是否属于影响集唯一确定，
集内占用没有指向集外的依赖边，因而「改写哪些」对应一个可检验的集合，
而不是人为设定的搜索半径，也不随改路算法的实现而变。
\emph{影响集作为划界对象}的证据是：走廊阻断下工序依赖图上的范围为空
而原排程已不可行（\(\PassMiss\)，第~\ref{sec:e1}~节）；预约影响集对外封闭（\(\PassClosed\)，
第~\ref{sec:e2}~节）；A 类扰动下预约影响集包含工序依赖图已纳入的全部占用（第~\ref{sec:e6}~节）；
边界差分试探的泄漏为 \ProbeLeakPairs。
\emph{该划界的净收益}由与 RA 的对照给出（第~\ref{sec:cheap}~节），
即相对「改写全部尚未结束的占用」减少的改动。

\item \textbf{给出与该划界相匹配的局部恢复，并分清各项结果各来自何处。}
在影响集内为原车辆改路、集外保持原占用；仍不可行时再扩大允许改写的范围。
修复步骤沿用既有做法，本项贡献在于通过对照区分各项结果的来源。
在同一套走廊路网、时间窗路由与无冲突校验下对照表明：
按工序依赖图划界时走廊阻断 \(\FeasRone\) 恢复可行，按预约影响集划界时 \(\FeasRtwo\) 恢复可行，
可行与否来自划界对象，而不是改路程序的优劣；
有界修复以毫秒级耗时恢复可行，但完工时间 \(C_{\max}\) 较全局重调度更差，
且放宽预算不出现反超（\(\EfiveCells\) 个算例--随机种子格，所测预算区间内无交叉点）；
与「改写全部尚未结束的占用」的对照相比，耗时与 \(C_{\max}\) 几乎不变，
只有改动的减少可归因于影响集（第~\ref{sec:cheap}~节）。
在四类扰动与两种划界的交叉中，工序依赖图上的范围仅在走廊阻断与通行变慢下为空，
预约影响集均非空，且在机械臂故障与车辆故障下包含工序依赖图上的范围。
\end{enumerate}

同一车间、同一预约表上，开工前的一体化排程用预约表构建无冲突方案~\cite{vanos2026exact,li2026framework,sun2023integrated}；本文在执行时刻用预约表定位失效范围，二者作用于不同阶段。

\subsection{篇章结构}

第~\ref{sec:related_work}~章按四类回顾相关工作：作业车间如何按工序或时间轴划定重调度范围、
铁路调度如何表达排他资源上的占用先后、无冲突多车路径如何做局部重规划，
以及考虑运输的柔性作业车间把边界划在何处，并据此说明本文的位置。
第~\ref{sec:problem}~章给出走廊预约设定与执行期恢复问题：
车间、走廊预约表与工序依赖图，两类扰动，以及恢复须满足的条件。
第~\ref{sec:strc}~章定义时空预约闭包及其结果集合预约影响集，说明其依赖边与对外封闭等性质
（图~\ref{fig:graphs}、\ref{fig:closure}）。
第~\ref{sec:repair}~章说明如何在影响集内为原车辆改路、集外保持原占用，
以及仍不可行时如何扩大允许改写的范围，并与全局重调度对照（图~\ref{fig:framework}）。
第~\ref{sec:experiments}~章报告可行性、算法耗时、占用改动比例与完工时间，
并用案例甘特图在工序时间轴上展示统计结果对应的改动形态。
第~\ref{sec:conclusion}~章总结全文，并讨论局限与后续工作。

